\begin{frame}{val 꺾임이 에포크를 정한다 (6.3절)}
  \begin{columns}[T]
    \begin{column}{0.55\textwidth}
      \centering
      \includegraphics[width=0.92\linewidth]{ch16_1_dl_curves.png}
      \vskip0.3em
      {\footnotesize 왼쪽: GBDT 곡선 / 오른쪽: CNN 에포크 곡선 —
      val 손실은 꺾임 포인트 이후 상승}
    \end{column}
    \begin{column}{0.44\textwidth}
      \begin{block}{6.3절의 원칙이 그대로}
        ``train은 \textbf{배운 것을}, val만 \textbf{재는 것을}
        한다''. 에포크를 늘리면 train 손실은 0으로 가지만
        val 손실은 어느 에포크 이후 \textbf{오름}(과적합,
        Ch09.3).
      \end{block}
      \vskip0.4em
      {\small 이 val \textbf{꺾임 포인트}가 에포크·학습률·
      정규화 강도를 정하는 근거. M3의 ``학습 곡선 1장''은
      이 꺾임을 보여야 한다.}
    \end{column}
  \end{columns}
\end{frame}
